% Related work — full draft v1 (2026-08-28), from paper/related_work_notes.md.
% Five threads; each paragraph ends with the differentiation sentence.

\paragraph{Resampling probes of reasoning traces.}
Truncating a trace and sampling continuations has become the standard
counterfactual tool for interpreting chain-of-thought. \citet{lanham2023}
introduced early-answering truncation curves to measure how much the stated
reasoning matters; \citet{bogdan2025anchors}
resample sentence-level alternatives to attribute outcome changes to
individual steps; \citet{bigelow2025forking} resample at every token to
locate ``forking tokens'' where the outcome distribution shifts abruptly,
and \citet{bigelow2026forkingfast} show that much apparent per-token
volatility in such estimates is sampling noise. Closest in shape to our
probe, \citet{merrill2026ponr} fix prefixes and resample continuations at
scale to localize where a (deceptive) outcome ``locks in,'' and
\citet{ballon2026} read answer distributions from truncated prefixes,
finding that commitment rises steadily with prefix length; and
\citet{wang2026re2}, as motivation for training models to abandon bad
prefixes, compare continuing truncated \emph{incorrect} traces against
re-solving from scratch --- the nearest published continue-versus-restart
comparison, without matched token budgets or per-anchor value curves. All
of
these estimate what a prefix \emph{leads to}; none of them ask what the
prefix is \emph{worth} relative to not having it. Our contribution is the
missing control: an empty-prefix restart baseline $R(C)$ at matched total
generated-token budget, which converts per-anchor solve rates into a value measurement and
reveals that most apparent within-trace transitions are budget artifacts.

\paragraph{Rollout value estimates in process supervision.}
Our estimator --- the fraction of budgeted continuations from a partial
solution that reach the correct answer --- is the workhorse of automatic
process supervision: Math-Shepherd's soft label \citep{wang2024mathshepherd}
is exactly this quantity (its main experiments binarize it to whether any
continuation succeeds), OmegaPRM \citep{luo2024omegaprm} binary-searches
it for the first error, Phi-4's pivotal token search \citep{abdin2024phi4}
mines sharp jumps in it for preference pairs, and \citet{setlur2025progress}
formalize its increments as ``progress'' rewards. This literature treats the
quantity as training signal and its unreliability as label noise to be cleaned
\citep{zhang2025prmlessons}. We repurpose the same estimator as a
\emph{measurement instrument} --- with anchors on the model's own trace,
interval censoring, attempt-pooling, and the restart control --- and treat
its intermediate values as signal: roughly a third of ambiguous states
remain intermediate after eight attempts (success 0.3--0.6), which
undermines
the single-crossing, monotone-value picture implicit in binary-search
labeling schemes \citep{luo2024omegaprm} and in snowball-error theories of
monotone degradation \citep{gan2025snowball}.

\paragraph{Aha moments and their skeptics.}
The claim that reinforcement-trained reasoners exhibit emergent ``aha
moments'' originates in DeepSeek-R1's training anecdotes
\citep{deepseekr1}. Skepticism so far has been correlational or
lexical: reflection keywords exist before any RL \citep{liu2025noaha},
many self-verification steps are causally inert \citep{zhao2025fakeaha,
kang2025firsttry}, and at trace scale, mid-reasoning strategy shifts are
rare and associated with \emph{lower} accuracy \citep{daliberti2026illusion}.
We supply the counterfactual, compute-matched version of this skepticism:
rather than classifying shifts in observed traces, we measure whether any
anchor of a trace carries value that fresh computation at equal total budget
cannot reproduce, and find such anchors in 1 of 178 cells. Sharp
within-trace transitions reported by entropy change-point analyses
\citep{xu2026entropy} and commitment-boundary probes \citep{scalena2026}
are consistent with our data --- but our restart control shows that the
states they mark are, almost always, reachable from scratch at the same
total cost.

\paragraph{Sequential versus parallel test-time compute.}
Aggregate comparisons of long chains against independent samples are now
common: revision models conditioned on complete prior answers
\citep{snell2024}, repeated-sampling coverage laws \citep{brown2024monkeys},
budget forcing \citep{muennighoff2025s1}, and matched-budget comparisons
finding either sequential \citep{sharma2025sequentialedge} or parallel
\citep{ghosal2025mirage, pgap2026} advantages, with context contamination
as one proposed mechanism for failed retries \citep{yang2026retrying} and
reduced exploration when conditioning on prior answers \citep{pgap2026}. These comparisons operate at benchmark
level and condition on \emph{completed} answers. Our instrument moves the
comparison inside the trajectory --- the model's own truncated prefix
against a fresh start at matched total tokens, per problem --- and yields a
quantity the aggregate view cannot express: a prefix's value in fresh-token
units. The answer, for our models, is predominantly compression rather than
reachability (the prefix wins all nine exactly-matched comparisons, with
two wins and two ties on boundary proxies; an 8{,}192-token restart
reaches threshold in 11/13), which is empirical evidence for the no-benefit branch of the dichotomy of
\citet{wolf2026theory}: when self-reflection fails to reliably localize
early errors, conditioning on past attempts offers no asymptotic benefit
over independent restarts. Our per-problem restart dose--response curves also
give the sequential-parallel debate a diagnostic reading: one model's
failures dissolve with budget (compute-starved), the other's do not
(capability-limited), echoing overthinking \citep{chen2025overthinking}
and non-monotone thinking-length effects \citep{zhou2026morethinking}.

\paragraph{Predicting outcomes from internal signals.}
A rapidly growing line reports that hidden states predict reasoning
outcomes: probes at intermediate-answer positions reach AUROC $>0.9$
\citep{zhang2025knowright}, first-step probes 0.79
\citep{yuan2026diagnostic}, pre-generation probes beat question-text
baselines \citep{lugoloobi2026}, question-only probes predict success before
generation \citep{cencerrado2025noanswer}, and confidence signals drive
early exit and trace selection in systems work \citep{fu2024certaindex,
fu2025deepconf}. With few exceptions these positives
carry no problem-difficulty control, and the exceptions control only
partially: \citet{yuan2026diagnostic} run their within-problem analysis on
full-trace probes (their early-window positive is uncontrolled, and their
within-problem effect sizes fall as low as $d{=}0.13$), while
\citet{lugoloobi2026} probe strictly before generation and themselves
conclude that activations encode model-specific \emph{difficulty}. Our
pre-registered, single-shot test supplies the missing control at the point
where it bites: frozen
summary features of the run's internal dynamics add no detectable AUROC
over a difficulty baseline on held-out problems at the pre-registered
forecast point, and a labeled post-hoc sweep finds no rescuing window
from 128 to 2{,}048 tokens --- consistent with reports that
correct and incorrect trajectories diverge late \citep{sun2026trajectories},
that trajectory geometry remains systematically coupled to difficulty once
length is residualized \citep{gjolbye2026}, and that difficulty-driven sample composition contaminates position-wise
probe curves, as \citet{david2025temporal} candidly note of their own
later-window decline --- their early-window positive carries no
question-only control, and it is the result we replicate exactly and then
reduce to chance within problem (Section~\ref{sec:results-prediction}). Existing critiques of correctness
probes attack two other axes: their signals resist causal use
\citep{yuan2026diagnostic}, and under contaminated prefixes they track
coherence rather than grounded correctness \citep{pair2026}. The
difficulty axis has remained open. Meanwhile, difficulty
itself is directly decodable from hidden representations
\citep{llmknowsdifficulty2025, geometrichardness2026}, which is precisely
the mechanism our dissection isolates.
